\documentclass[12pt]{article}

%% Packages
\usepackage[margin=1.25in]{geometry}
\usepackage{setspace}
\usepackage{hyperref}
\usepackage{natbib}
\usepackage{amsmath}
\usepackage{amssymb}
\usepackage{microtype}
\usepackage{booktabs}
\usepackage{xcolor}
\usepackage{parskip}

\hypersetup{
  colorlinks=true,
  linkcolor=blue!60!black,
  citecolor=blue!60!black,
  urlcolor=blue!60!black
}

\onehalfspacing

\title{\textbf{Warehouse Certificates, Real Bills, and Fiduciary Money:\\
The Smith--Ricardo Doctrine from Ricardo's Ingot Plan\\
to Bretton Woods and Beyond}}

\author{Thomas J.\ Sargent}
\date{March 2026}

\begin{document}

\maketitle

\begin{abstract}
\noindent
David Ricardo's 1816 ``Ingot Plan'' proposed that Britain replace
circulating gold coin with paper bank notes backed not by 100-percent gold
reserves but by holdings of short-term, self-liquidating commercial loans
(``real bills'' in Adam Smith's sense) together with a modest gold reserve
sufficient only for international settlements.  Ludwig von Mises in
\textit{Human Action} (1949) called this combination of Smith's banking
theory with Ricardo's monetary proposal the ``Smith--Ricardo doctrine''
and identified it as the intellectual parent of three successive
international monetary regimes: the pre-1914 Gold Exchange Standard, the
British return to gold in the 1920s, and the Bretton Woods system
engineered by John Maynard Keynes and Harry Dexter White that governed IMF
policies from 1944 until Richard Nixon closed the gold window in 1971.
Irving Fisher and Keynes each argued, with somewhat different emphases,
that a well-managed fiduciary currency system could strictly dominate a
100-percent metallic standard on both efficiency and price-level-stability
grounds---a proposition that might be called the Fisher--Keynes Theorem of
monetary superiority.  Milton Friedman acknowledged the efficiency claim
but was deeply ambivalent about stability, emphasizing the discretionary
risks inherent in any managed currency.  The 1982 paper by Thomas Sargent
and Neil Wallace reconsidered the real-bills prescription within an
overlapping-generations framework and found it Pareto-superior to
quantity-theory restrictions under certain conditions, while David Laidler
immediately cautioned that their formal doctrine differed substantially
from the 19th-century original.  This essay reconstructs that nexus of
ideas and identifies the unresolved tensions that run through it.
\end{abstract}

\tableofcontents
\newpage

%% ------------------------------------------------------------------
\section{Introduction}
\label{sec:intro}
%% ------------------------------------------------------------------

There is a connected set of monetary doctrines, originating in
18th-century British political economy, that has shaped every major
international monetary system of the past two centuries.  At its core
stands a deceptively simple claim: that a country can do better---both
economically and in terms of price stability---by circulating a managed
paper currency backed by productive-activity loans than by hoarding the
gold and silver that such currency replaces.  The logic is compelling.
A 100-percent gold reserve system ties up a valuable commodity in
vault storage, sacrificing the real resources that gold could otherwise
mobilize.  If a trusted issuer can replace gold coin with paper
promises that circulate at face value, the gold can be exported to buy
real imports, freeing up the ``dead stock of gold and silver'' (in
Adam Smith's phrase) for productive use.

Adam Smith laid out the banking theory underlying this claim in
\textit{The Wealth of Nations} (1776).  David Ricardo translated it
into a concrete institutional proposal in his \textit{Proposals for an
Economical and Secure Currency} (1816).  John Maynard Keynes formalized
the gold-exchange standard in his \textit{Indian Currency and Finance}
(1913) and explicitly credited Ricardo with its theoretical foundations.
Irving Fisher argued in \textit{The Purchasing Power of Money} (1911)
and elsewhere that a properly managed fiduciary system would stabilize
the price level better than any metallic standard, because the money
supply could be adjusted to offset shocks rather than varying
passively with gold discoveries and trade flows.

Ludwig von Mises, in \textit{Human Action} (1949), brought these
threads together under the label ``the Smith--Ricardo doctrine'' and
argued that the doctrine was a ``fallacy'' whose misapplication had
generated costly monetary instabilities from the 19th century through
the collapse of Bretton Woods.  This essay traces the doctrine's
intellectual genealogy, examines the major monetary systems it
informed, surveys the competing modern assessments of its validity,
and identifies the irreducible tensions between efficiency and
stability that lie at its heart.

%% ------------------------------------------------------------------
\section{Adam Smith's Real Bills Doctrine}
\label{sec:smith}
%% ------------------------------------------------------------------

\subsection{The Argument in \textit{The Wealth of Nations}}

Book II, Chapter II of \textit{The Wealth of Nations} (1776) contains
the essential statement of what Lloyd Mints \citeyearpar{Mints1945}
would later term the ``real bills doctrine.''  Smith's argument begins
with the observation that a country's circulating medium is itself a
part of its fixed capital: gold and silver coins are productive
assets only insofar as they facilitate exchange.  Any quantity of coin
above the minimum required for the ``wheel of circulation'' is,
from the social standpoint, a waste.

Banks can reduce this waste by substituting paper for metal.  When a
bank discounts a bill of exchange---a short-term instrument
representing goods already produced and in transit to market---it
issues bank notes that circulate in place of coin.  Because the bill
is self-liquidating (it is repaid when the underlying goods are sold,
typically within 60 to 90 days), the bank notes are backed by a real
claim on output.  Smith writes:

\begin{quote}
The gold and silver money which circulates in any country may very
properly be compared to a highway, which, while it circulates and
carries to market all the grass and corn of the country, produces
itself not a single pile of either.  The judicious operations of
banking, by substituting paper in the room of a great part of this
gold and silver, enables the country to convert, as it were, a great
part of its highways into good pastures and corn fields, and thereby
to increase very considerably the annual produce of its land and
labour. \citep[Book~II, Ch.~II]{Smith1776}
\end{quote}

Two features of Smith's argument deserve emphasis.  First, the
safety of the bank does not require a 100-percent gold reserve; it
requires only that the bills discounted be genuine commercial claims
on real output, not accommodation paper issued to finance consumption
or speculation.  Second, the quantity of notes issued will be
automatically self-regulating: notes return to the bank for
redemption precisely when the goods they financed are sold.  The
note supply thus tracks the ``needs of trade'' without deliberate
central management.

\subsection{The Three Flaws: Thornton's Critique}

The doctrine's critics identified three fundamental weaknesses.
Henry Thornton, in \textit{An Enquiry into the Nature and Effects of
the Paper Credit of Great Britain} (1802), first articulated what
would become the standard critique.

\textbf{First}, the doctrine confuses real with nominal values.  Notes
are backed by the nominal value---price times quantity---of commercial
bills, not by their real value alone.  If price levels rise, the
nominal value of bills rises, note issue expands to accommodate the
higher price level, and a self-reinforcing inflationary spiral is
possible whenever the market (loan) rate of interest is held below the
natural (profit) rate of capital \citep{Thornton1802}.

\textbf{Second}, the distinction between ``genuine'' commercial bills
and ``speculative'' or ``accommodation'' paper is practically
impossible to enforce.  A banker cannot reliably distinguish a bill
backed by goods in transit from one backed by a forward bet on
commodity prices; the paper looks identical.

\textbf{Third}, and most consequentially, the doctrine is
pro-cyclical.  When activity expands, bill volumes rise, note issue
expands, and prices are pushed further upward.  When activity
contracts, the reverse occurs: note issue shrinks precisely when
a stabilizing expansion of credit might arrest the downturn.  Modern
stabilization theory calls for counter-cyclical money, but the real
bills doctrine delivers pro-cyclical money.

These criticisms did not immediately dislodge the doctrine.  It
survived, transmuted and rehabilitated, through the 19th-century
currency--banking debate in Britain and re-emerged as a guiding
principle of the Federal Reserve Act of 1913
\citep{Fetter1965, Humphrey2001}.

%% ------------------------------------------------------------------
\section{Ricardo's Ingot Plan: Paper Warehouse Certificates for Gold}
\label{sec:ricardo}
%% ------------------------------------------------------------------

\subsection{The 1816 Proposal}

David Ricardo's \textit{Proposals for an Economical and Secure
Currency: With Observations on the Profits of the Bank of England}
(1816) represented the most rigorous translation of Smith's banking
theory into a concrete monetary architecture.  The proposal emerged
directly from the Bullionist Controversy (roughly 1797--1821), the
great British debate about whether the wartime suspension of
convertibility had caused the depreciation of the pound relative
to gold.

Ricardo was the leading ``Bullionist''---he argued, correctly in the
event, that the Restriction Period (1797--1821) had allowed the Bank
of England to over-issue notes, depreciated the pound, and raised
prices.  But his remedy was not a simple return to full gold
circulation.  That would waste resources.  Instead, Ricardo proposed
what he called the ``Ingot Plan'':

\begin{enumerate}
  \item Bank notes would remain the sole circulating medium; gold
        coin would be replaced by paper.
  \item The Bank of England would be obliged to redeem its notes in
        \emph{gold bullion bars} (ingots), not in coins, at a fixed
        price.  The minimum redemption unit would be set large
        enough---400 ounces in the original proposal---to ensure that
        gold flowed to the Bank only for the purpose of international
        settlement, not domestic face-to-face transactions.
  \item The Bank would hold a modest gold reserve sufficient for
        settling international imbalances.  Domestic circulation
        would be paper throughout.
  \item The Bank would expand and contract its note issue by
        discounting and un-discounting real commercial bills of
        exchange, thereby backing currency with productive activity
        rather than idle gold reserves.
\end{enumerate}

As Keynes later wrote, Ricardo ``laid it down that a currency is in
its most perfect state when it consists of a cheap material, but
having an equal value with the gold it professes to represent; and he
suggested that convertibility for the purposes of the foreign
exchanges should be ensured by the tendering on demand of gold
\emph{bars} (not coin) in exchange for notes, so that gold might be
available for purposes of export only, and would be prevented from
entering into the internal circulation of the country''
\citep[Ch.~II]{Keynes1913}.

The warehouse certificate metaphor captures the economic logic
precisely.  A depositor of gold at the Bank receives a paper
certificate acknowledging the deposit; that certificate circulates
as money.  But Ricardo's certificates are not purely warehouse
receipts: the Bank also issues certificates against commercial
loans.  The gold reserve is held fractionally, not fully.  The
``certificate'' is partly backed by gold and partly by the
productive credit of merchants---precisely Adam Smith's vision.

\subsection{The Bullionist Controversy and Its Resolution}

The Bullion Committee Report of 1810, drafted largely on the
basis of Ricardo's testimony and arguments, endorsed the Bullionist
position: the pound had depreciated because of excess note issue
by the Bank \citep{BullionReport1810}.  Its recommendation for
an early resumption of convertibility was initially rejected by
Parliament, embarrassing the Bank, but eventually vindicated by
the return to gold in 1819 at the pre-war parity.

The Ingot Plan itself was never implemented exactly, but it
provided the intellectual template for subsequent monetary
architecture.  The Bank Charter Act of 1844 (``Peel's Act'')
moved in the opposite direction---toward 100-percent marginal
gold backing for any note issue above a fixed fiduciary
issue---but the gold exchange standard that gradually emerged
in practice among Britain's trading partners was far closer
to Ricardo's vision than to Peel's.

%% ------------------------------------------------------------------
\section{The Mises Critique: A Doctrine and Its Fallacy}
\label{sec:mises}
%% ------------------------------------------------------------------

Ludwig von Mises, in \textit{Human Action: A Treatise on Economics}
(1949), gave the tradition its sharpest critical label.  He
identified a ``Smith--Ricardo doctrine'' that he traced from Smith's
banking chapters through Ricardo's Ingot Plan to the 20th-century
monetary systems he regarded as fatally flawed.

Mises's argument was twofold.  First, the doctrine was technically
defective: the real bills criterion cannot stabilize the money
supply because it allows note issue to expand endlessly so long as
bankers continue to lend against commercial paper whose nominal
value inflates with prices \citep{Mises1949}.  This repeats
Thornton's first objection in Austrian garb.

Second, and more fundamental for Mises, the doctrine provided cover
for central bank discretion and inflationary finance.  Once it is
accepted that a trusted central bank can manage a fractional-reserve
currency better than an automatic gold standard, the door is open
to successive debasements justified by appeals to productive
necessity.  The Smith--Ricardo doctrine was, for Mises, not merely
an intellectual error but an institutional temptation---a
``fallacy'' that historically supported each successive erosion of
gold discipline.

Mises traced this institutional path explicitly:

\begin{enumerate}
  \item  The \emph{Gold Exchange Standard}, described by Keynes in
         \textit{Indian Currency and Finance} (1913) and popularized
         in the late 19th century especially through the Indian and
         Dutch currency reforms, allowed peripheral countries to hold
         not gold itself but claims on a gold-standard country (sterling)
         as monetary reserves.  The Smith--Ricardo logic justified
         the arrangement: a currency is ``in its most perfect state''
         when it consists of a cheap material equal in value to gold,
         without gold itself circulating domestically.
  \item  The \emph{British return to gold in 1925} at the old parity
         of \$4.86 per pound was arranged on gold-exchange-standard
         principles, with Britain holding significant foreign-exchange
         reserves alongside a diminished gold hoard.  Winston Churchill,
         as Chancellor of the Exchequer, implemented the Gold Standard
         Act of 1925 over the famous objections of Keynes, who argued
         in \textit{The Economic Consequences of Mr.\ Churchill} (1925)
         that the overvalued parity would impose severe deflation.
  \item  The \emph{Bretton Woods system} (1944--1971) extended the
         gold-exchange logic globally.  Other currencies fixed to
         the U.S.\ dollar; only the dollar was convertible into gold,
         and only for foreign central banks, at the fixed rate of
         \$35 per troy ounce.  The IMF enforced exchange-rate
         discipline within this dollar-gold pyramid.
\end{enumerate}

For Mises, each step in this sequence represented a further
removal from the discipline of a full gold standard, and each was
legitimized by the Smith--Ricardo argument that managed paper
is superior to metallic money---a claim he regarded as the root
of 20th-century monetary disorder \citep{Mises1949,
Salerno1982}.

%% ------------------------------------------------------------------
\section{The Gold Exchange Standard and Keynes}
\label{sec:ges}
%% ------------------------------------------------------------------

\subsection{Keynes's 1913 Codification}

John Maynard Keynes's \textit{Indian Currency and Finance} (1913)
is the most authoritative early codification of the gold exchange
standard as a monetary system.  Writing about the Indian rupee
(fixed to sterling in 1893), Keynes described the theoretical
advantages that Ricardo had identified nearly a century earlier:

\begin{quote}
The Gold-Exchange Standard arises out of the discovery that, so
long as gold is available for payments of international
indebtedness at an approximately constant rate in terms of the
national currency, it is a matter of comparative indifference
whether it actually forms the national currency\ldots\ The
Gold-Exchange Standard may be said to exist when gold does not
circulate in a country to an appreciable extent, when the local
currency is not necessarily redeemable in gold, but when the
Government or Central Bank makes arrangements for the provision
of foreign remittances in gold at a fixed maximum rate in terms
of the local currency, the reserves necessary to provide these
remittances being kept to a considerable extent abroad.
\citep[Ch.~II]{Keynes1913}
\end{quote}

Keynes was explicit that this system's ``theoretical advantages
were first set forth by Ricardo.''  He noted that among its
virtues were (a) economy in the use of gold---the resource cost
of maintaining a full gold circulation was equivalent to tying up
capital in idle metal---and (b) stability, since a managed reserve
could be held at whatever level was consistent with exchange
stability rather than being thrown about by arbitrary gold flows.

\subsection{Britain's Return to Gold, 1925}

The Cunliffe Committee (1918--1919) recommended that Britain
return to gold at the pre-war parity after the war.  When
Winston Churchill implemented this in April 1925, the arrangement
was essentially a gold bullion standard rather than a gold specie
standard: the Bank of England held gold bars, not coins, and
would exchange notes for bars only in large minimum lots.
Domestic circulation was entirely in paper.  This was almost
exactly Ricardo's Ingot Plan realized 109 years after the fact.

Keynes objected strenuously \citep{Keynes1925}, not to the
principle of a gold-backed managed currency, but to the choice
of the pre-war parity.  He argued that the pound was overvalued
by approximately 10 percent at \$4.86, and that restoring the
old parity would require severe domestic deflation to restore
competitiveness---deflation that would fall disproportionately on
wages and employment given labor market rigidities.  He was proved
correct: the overvalued pound contributed to chronic British
stagnation in the late 1920s and forced the final departure from
gold in September 1931 \citep{Eichengreen1992}.

\subsection{The Interwar Period and the Collapse}

The interwar gold exchange standard, reconstructed at Genoa (1922)
under Anglo-American auspices, was a pale shadow of the pre-1914
classical gold standard \citep{Eichengreen1992, Nurkse1944}.
Several structural problems undermined it:

\begin{enumerate}
  \item \textbf{Reserve concentration}:  Gold was heavily
        concentrated in the United States and France.  The price-specie
        flow mechanism required that gold-rich countries inflate while
        gold-poor countries deflate, but both the Federal Reserve and
        the Banque de France ``sterilized'' their inflows to avoid
        domestic inflation.
  \item \textbf{Short-term capital flows}:  The sterling-dollar
        exchange was subject to speculative attack in a way the
        pre-war system, backed by deeper credibility, had largely
        escaped.
  \item \textbf{Real bills policy at the Fed}:  Under the influence
        of the real bills doctrine (particularly the ``Direct Pressure''
        initiative of Federal Reserve Board member Adolph C.\ Miller),
        the Fed contracted credit in 1929--1931 precisely when
        expansion was needed, contributing to the Great Contraction
        analyzed by \citet{FriedmanSchwartz1963} and later
        \citet{HumphreyTimberlake2019}.
\end{enumerate}

%% ------------------------------------------------------------------
\section{Bretton Woods: Keynes, White, and the Dollar-Gold System}
\label{sec:bretton-woods}
%% ------------------------------------------------------------------

The Bretton Woods Conference of July 1944 created the post-World
War II monetary order that governed international finance until
Nixon's announcement of August 15, 1971.  The intellectual
architects were John Maynard Keynes (representing Britain) and
Harry Dexter White (representing the United States), and their
design bore the clear imprint of the Smith--Ricardo tradition.

Keynes's draft plan for an \textit{International Clearing Union}
(the ``Keynes Plan'') proposed a supranational institution with
its own unit of account (the ``bancor''), whose members would
hold bancor balances rather than gold.  Trade imbalances would
be settled in bancor, with symmetric adjustment obligations on
both surplus and deficit countries.  Gold would play no active
circulating role; it would be a one-way valve used by countries
to acquire bancors but not the reverse \citep{Keynes1943}.

White's draft plan was less radical but structurally similar
in its drift away from a full gold standard.  The compromise
agreed at Bretton Woods established:

\begin{enumerate}
  \item  Fixed exchange rates for all member currencies denominated
         in U.S.\ dollars.
  \item  Dollar convertibility into gold at \$35 per ounce, but
         only for foreign monetary authorities, not private parties.
  \item  An International Monetary Fund (IMF) to provide short-term
         balance-of-payments financing and supervise exchange
         rate adjustments.
  \item  Capital controls permitted (indeed, expected) as part of
         the adjustment mechanism.
\end{enumerate}

The arrangement was again a gold exchange standard in Ricardo's
sense: gold did not circulate domestically in any country; it served
only as the ultimate reserve behind the dollar, behind which all
other currencies stood.  The dollar was, in effect, a warehouse
certificate for gold, and other currencies were warehouse certificates
for dollars.

The system worked tolerably well through the 1950s and
into the 1960s, but its design contained a structural paradox
identified by Robert Triffin in 1960 (the ``Triffin dilemma''):
the world's demand for dollar reserves could only be satisfied
by persistent U.S.\ balance-of-payments deficits, but those
deficits would eventually erode confidence in dollar
convertibility \citep{Triffin1960}.  As U.S.\ gold reserves
fell from about 574 million troy ounces in 1949 to 296 million
by 1971, and as foreign dollar claims accumulated, the system
became increasingly fragile.  The final crisis came with
France's persistent demands for gold conversion of its dollar
holdings (de Gaulle's ``assault on Fort Knox'') and the
fiscal excesses of the Vietnam War era.  Nixon closed the
gold window on August 15, 1971, and the Bretton Woods
system---the culmination of the Smith--Ricardo institutional
lineage---ended \citep{Eichengreen1996}.

%% ------------------------------------------------------------------
\section{The Fisher--Keynes Theorem: Fiduciary Money Dominates}
\label{sec:fisher-keynes}
%% ------------------------------------------------------------------

\subsection{Fisher's Stability Case}

Irving Fisher in \textit{The Purchasing Power of Money} (1911)
and \textit{Stabilizing the Dollar} (1920) went further than
Smith and Ricardo in arguing for the superiority of managed paper.
He proposed the ``compensated dollar'' plan: adjust the gold
content of the dollar periodically to stabilize the price level.
Under a 100-percent gold standard, price level movements are
governed by the ratio of gold production to real output growth.
Under Fisher's managed standard, a central authority adjusts the
gold exchange rate to keep the price level stable---achieving
both the resource cost savings of fiduciary money and the
price-level stability that a metallic standard can only deliver
accidentally.

Fisher's implicit ``theorem'' (he did not state it formally)
is that a managed fiduciary currency weakly dominates a 100-percent
metallic standard:
\begin{itemize}
  \item It saves the resource cost of the gold reserve
        (efficiency advantage);
  \item It achieves price stability more reliably, since money
        can be contracted or expanded to offset shocks rather
        than varying with the accidents of gold mining
        (stability advantage).
\end{itemize}

Both advantages hold if---and this is the critical assumption---the
monetary authority is competent and honest.  Fisher was not
na\"ive about this; his 100-percent reserve banking proposal
of 1935 was precisely an attempt to discipline the banking system
by separating the money-creation function from the
lending function \citep{Fisher1935}.

\subsection{Keynes's Parallel Argument}

Keynes made a parallel argument with characteristic breadth.
In \textit{A Tract on Monetary Reform} (1923) he wrote:

\begin{quote}
In truth, the gold standard is already a barbarous relic.  All of
us, from the Governor of the Bank of England downwards, are now
primarily interested in preserving the stability of business,
prices, and employment, and are not likely, when the choice is
forced on us, deliberately to sacrifice these to the outworn
dogma, which had its value once, of \pounds 3 17s.\ 10\textonehalf d.\ per ounce. 
\citep{Keynes1923}
\end{quote}

The argument was not merely rhetorical.  In \textit{A Treatise on
Money} (1930) and \textit{The General Theory} (1936) Keynes
developed a more systematic case: a gold standard imposes
a deflationary bias in a world where wages are sticky downward,
because trade imbalances must be corrected by price and wage
adjustment borne by deficit countries, not by symmetric expansion
in surplus countries.  A fiduciary system with a flexible exchange
rate---or an internationally coordinated system with symmetric
adjustment obligations---could in principle eliminate this
deflationary bias while preserving monetary discipline.

The efficiency argument (saving the resource cost of idle gold
reserves) also appears in Keynes, most explicitly in the Bretton
Woods negotiations, where he estimated the global welfare gain
from mobilizing that gold for productive investment.

\subsection{Is It a Theorem?}

The Fisher--Keynes claim is heuristically compelling but not a
formal theorem in the sense of a precisely stated proposition
with clearly articulated hypotheses and a proof.  The nearest
formal statement in modern monetary theory is probably the
following \citep[cf.][]{FriedmanSchwartz1963, Lucas1992}:

\begin{quote}
\textit{If the monetary authority can credibly commit to a
price-level or inflation target, then a fiduciary currency
weakly dominates a commodity-money standard: it provides the
same (or better) price-level stability at lower resource cost,
because the commodity reserve is eliminated and the currency
supply can be adjusted in response to demand shocks.}
\end{quote}

The key conditional---``if the monetary authority can credibly
commit''---is where all the institutional difficulty lies.  A
commodity standard provides credibility through a physical
constraint; a fiduciary standard must generate credibility
from institutions, reputations, and constitutional arrangements
that are harder to establish and easier to abandon
\citep{Bordo1990, BarroGordon1983}.

The claim is also not obviously true in all environments.  In an
economy where the fiscal authority is prone to seigniorage
extraction---financing deficits by money creation---a
hard metallic constraint may dominate a fiduciary system not
because of its efficiency properties but because of its
commitment properties \citep{Sargent1982unpleasant}.  This is
the core of Mises's critique: in a world of governments
susceptible to inflationary pressure, the Fisher--Keynes comparison
of a \emph{well-managed} fiduciary system against a metallic
standard is misleading, because the metallic standard is the
relevant commitment device precisely in environments where
``well-managed'' cannot be taken for granted.

%% ------------------------------------------------------------------
\section{Friedman's Ambivalence}
\label{sec:friedman}
%% ------------------------------------------------------------------

Milton Friedman's position on the Fisher--Keynes claim is genuinely
complex and evolved over his career.  His ambivalence reflects
the irreducible tension between the efficiency gains from fiduciary
money and the commitment problems that fiduciary money creates.

On the efficiency side, Friedman was an early and consistent
advocate of a pure fiat currency eliminating the resource cost
of commodity money.  In \textit{A Program for Monetary Stability}
(1960) he proposed a fixed-rate monetary growth rule---what
became known as the ``k-percent rule''---as the best practical
substitute for the automatic discipline of a gold standard
\citep{Friedman1960}.  The proposal accepted the Fisher--Keynes
efficiency argument: there is no reason to tie up real resources
in gold when paper can do the same monetary work.

On the stability side, however, Friedman was deeply skeptical
of managed discretionary monetary systems---including exactly
the managed gold exchange standards that Ricardo, Keynes, and
Fisher had advocated.  His historical case study was damning.
In \textit{A Monetary History of the United States, 1867--1960}
(1963), co-authored with Anna J.\ Schwartz, Friedman showed that
the Federal Reserve---operating under real-bills principles in
the 1930s---had caused the Great Contraction through errors of
commission (passive acquiescence to bank failures) and omission
(failure to expand the money supply) \citep{FriedmanSchwartz1963}.

\subsection{The Optimum Quantity of Money}

Friedman's analysis of the \textit{optimum quantity of money}
(1969) provides the clearest statement of the efficiency dimension
of his thought.  The social opportunity cost of holding money is
the nominal interest rate; the marginal cost of producing fiat
money is approximately zero.  Therefore the optimal monetary
policy requires driving the nominal interest rate to zero through
mild deflation at the rate of the real interest rate
\citep{Friedman1969}.  A metallic standard cannot achieve this
optimum without deflation equal to the real return on capital,
which is costly in a world of debt-deflation dynamics; a fiat
standard can achieve it by choice of the monetary growth rate.

\subsection{The k-Percent Rule as a Middle Ground}

Friedman's k-percent rule was an attempt to capture the
efficiency advantages of fiat money while avoiding the discretion
risks of managed money.  The rule was constitutional rather than
merely advisory: Friedman wanted the money growth rate fixed by
law, not by the judgment of central bankers.  In this respect
the rule is a substitute for the gold standard's automatic
discipline, not an endorsement of fine-tuned management.

His famous statement that ``the gold standard works
well under only those circumstances in which it is not needed''
captures his position: a gold standard is a superfluous constraint
when the fiscal and monetary authorities are already disciplined,
and an ineffective constraint when they are not, since a
determined government can always abandon convertibility
\citep{Friedman1962}.  Yet Friedman was also skeptical that
central bankers could be trusted to manage a fiat currency
wisely---hence the k-percent rule as a rule-based alternative
to both.

Late in life, Friedman modified even this position, suggesting
that monetary targeting had become impractical as financial
innovation blurred the boundaries of monetary aggregates, and
that inflation targeting by independent central banks might be
the best available approximation to his rule-based ideal
\citep{Friedman2003}.  His ambivalence about the Fisher--Keynes
claim thus never fully resolved: he accepted the efficiency
argument in principle while doubting that any real institution
could reliably deliver the stability benefits in practice.

%% ------------------------------------------------------------------
\section{Sargent--Wallace (1982) and the Modern Reassessment}
\label{sec:sargent-wallace}
%% ------------------------------------------------------------------

In 1982 Thomas Sargent and Neil Wallace published ``The Real-Bills
Doctrine versus the Quantity Theory: A Reconsideration'' in the
\textit{Journal of Political Economy} \citep{SargentWallace1982}.
Working within an overlapping-generations (OLG) model, they showed
that the traditional quantity-theory prescription---restrict private
intermediation to control the money supply---is not Pareto optimal,
while the real-bills prescription---allow unfettered private
intermediation or central bank operations designed to replicate such
intermediation---achieves the Pareto optimal allocation.

The intuition is that in an OLG economy, money serves not only as
a medium of exchange but as a store of value.  Restricting the
supply of money-substitutes (safe, short-term financial assets)
forces agents into an inferior portfolio, reducing welfare.
Allowing banks to issue notes against commercial loans---the
real-bills prescription---expands the supply of safe assets to
its efficient level.

Sargent and Wallace were careful to note that their result depends
on specific modeling assumptions, including the absence of aggregate
demand shocks that might make monetary control desirable for
stabilization purposes.  But the formal result was striking: under
those assumptions, the real bills prescription is superior to
quantity-theory restrictions.

Monetary historian David Laidler \citeyearpar{Laidler1984} immediately
published a comment noting that Sargent and Wallace's formal doctrine
differed substantially from the historical 19th-century real bills
doctrine.  The historical doctrine was about preventing over-issue
by tying notes to real commercial activity; Sargent and Wallace's
result was about achieving Pareto optimality through unrestricted
intermediation.  These are not the same claim, and the historical
doctrine's critics---Thornton, Fisher, Friedman---were targeting
the former, not the latter.

The exchange illustrates a general methodological point: formal
models can rehabilitate the \textit{intuition} behind an old
doctrine while differing materially from its historical content
and institutional applications.  Both the real bills doctrine's
historical defenders and its defenders in modern theory are pointing
at the efficiency gains from allowing the private financial system
to create liquid claims on real activity.  But the stability risks
that the historical critics identified are not captured in OLG
models without aggregate uncertainty, and it is precisely those
stability risks that determine the welfare comparison in practice.

%% ------------------------------------------------------------------
\section{The Unresolved Tensions}
\label{sec:tensions}
%% ------------------------------------------------------------------

The Smith--Ricardo doctrine traverses a remarkable landscape of
monetary thought.  Several fundamental tensions run through it that
remain unresolved.

\subsection{Efficiency versus Credibility}

The efficiency case for fiduciary money is robust: eliminating gold
from domestic circulation frees up real resources, and a managed
currency can in principle achieve price stability at lower deadweight
loss than an automatic metallic standard.  But the efficiency gains
are only realized if the monetary authority is credible---if agents
believe it will not exploit the liberty to print money for fiscal
purposes.  This credibility is harder to establish for a fiat
currency than for a metallic standard, because the metallic
obligation is physical while the fiat obligation is institutional.

\citet{BarroGordon1983} formalized this tension as the time-inconsistency
problem: a benevolent government that can issue fiat money at zero
cost faces a systematic incentive to surprise agents with inflation
to exploit the short-run Phillips-curve tradeoff.  Rational agents
anticipate this and require compensation through higher nominal wages
and interest rates, producing an inflationary equilibrium with no
employment gain.  The gold standard's discipline prevents this
outcome by eliminating the temptation.

\subsection{Stability versus Flexibility}

The deflationary bias of a metallic standard in a world of sticky
wages is the Keynesian case for fiduciary money.  This case was
spectacularly confirmed by the Great Depression: the gold standard
transmitted deflationary shocks internationally and constrained
central banks from offsetting them
\citep{BernankeJames1991, Eichengreen1992}.  Countries that
abandoned gold early (Britain in 1931, the United States in 1933)
recovered earlier.

But fiduciary flexibility without discipline produced the Great
Inflation of the 1970s: after Nixon closed the gold window, the
major economies experienced a decade of stagflation that the
Bretton Woods architects had not foreseen.  The discipline that
Mises saw as the gold standard's chief virtue proved to be the
Bretton Woods system's chief weakness.  Achieving both flexibility
and discipline simultaneously---what modern inflation-targeting
regimes attempt---requires institutional arrangements whose
permanence neither Smith, Ricardo, Fisher, nor Keynes had reason
to take for granted.

\subsection{The Real Bills Problem}

The real bills doctrine embedded in Ricardo's plan and in the
Federal Reserve Act of 1913 was precisely the mechanism that,
under the Federal Reserve's literal interpretation in the 1920s
and 1930s, generated the pro-cyclical contraction
\citet{HumphreyTimberlake2019} identify as a proximate cause of
the Great Depression.  The doctrine's recommendation to restrict
the discount window to ``real'' as opposed to ``speculative''
bills---fine in theory, impossible to implement in practice---led the
Fed to contract credit exactly when expansion was required.

The irony was total: a doctrine designed to ensure monetary
adequacy by tying note supply to productive activity ended up
producing monetary stringency by applying a distinction that was
doctrinally clear but operationally meaningless.  As Adolph
C.\ Miller's ``Direct Pressure'' letter of 1929 showed, the doctrine
was not merely irrelevant but actively harmful when wielded
by administrators who took its productive--speculative distinction
seriously \citep{HumphreyTimberlake2019}.

\subsection{International Architecture}

At the international level, the Smith--Ricardo doctrine generates
a different tension.  A gold exchange standard economizes on gold
but creates a hierarchy of credibility: peripheral countries hold
claims on center-country currencies rather than gold itself.  This
architecture is efficient only so long as the center country
maintains its gold commitment.  When that commitment becomes
incredible---as the Triffin dilemma explains for Bretton
Woods---the entire system is vulnerable to a sudden transition to
a new regime.  The 1971 Nixon shock was not a smooth policy
adjustment; it was a credibility crisis whose disorderly resolution
produced a decade of currency volatility \citep{Volcker1992}.

The modern post-Bretton Woods system of floating exchange rates and
inflation-targeting central banks is, in effect, the logical
endpoint of the Smith--Ricardo trajectory: all gold constraint has
been removed, and monetary discipline must be generated entirely
by institutional credibility.  Whether this represents the
Fisher--Keynes efficiency frontier or the Misesian nightmare of
unconstrained fiduciary money depends, as always, on the
quality of the governing institutions.

%% ------------------------------------------------------------------
\section{Conclusion}
\label{sec:conclusion}
%% ------------------------------------------------------------------

The Smith--Ricardo doctrine is one of the most consequential ideas
in the history of monetary policy.  Its core insight---that a trusted
institution can serve the monetary functions of gold at a fraction
of the resource cost---is correct.  Adam Smith saw it in the
operations of Scottish banks discounting commercial paper.  Ricardo
formalized it in the Ingot Plan.  Keynes extended it to international
monetary architecture.  Fisher argued it implied a superiority of
managed fiduciary money over automatic metallic standards.

Two centuries of experience have simultaneously confirmed the
efficiency insight and complicated the stability claim.  The Great
Depression demonstrated the deflationary dangers of an automatic
gold standard in a world of sticky wages and financial fragility.
The Great Inflation of the 1970s demonstrated the inflationary
dangers of a fiduciary system without adequate institutional
discipline.  The Fisher--Keynes argument that a well-managed
fiduciary system dominates a metallic standard is probably best
understood as a claim about the frontier of attainable outcomes:
if institutions are adequate, fiduciary money achieves both
greater efficiency and greater stability.  Whether institutions
are in fact adequate is an empirical question whose answer varies
by country, era, and constitutional arrangement.

Friedman's ambivalence is, in retrospect, the most honest position.
The efficiency claim is strong; the stability claim is conditional
on institutional quality; the historical record is mixed.  The
Smith--Ricardo doctrine illuminates the efficiency frontier of
monetary design.  It does not resolve the political economy problem
of how to get there and stay there.

The modern global monetary order---floating rates, independent
central banks, inflation targets---is the inheritor of the
Smith--Ricardo doctrine in its most ambitious form: a pure fiduciary
system globally, with gold playing no formal role.  Its performance
over the past quarter-century in the advanced economies has been
generally better than either the Great Depression gold standard or
the 1970s managed float would suggest.  Whether (in Mises's terms)
the doctrine ``fallacy'' has finally been tamed by institutional
design, or whether the Smith--Ricardo tensions are merely in
abeyance awaiting the next fiscal or financial shock, remains
to be seen.

%% ------------------------------------------------------------------
\bigskip
\bibliographystyle{plainnat}

\begin{thebibliography}{99}

\bibitem[Barro and Gordon(1983)]{BarroGordon1983}
Barro, R.~J. and D.~B. Gordon (1983).
\newblock Rules, discretion and reputation in a model of monetary policy.
\newblock \textit{Journal of Monetary Economics} 12(1), 101--121.

\bibitem[Bernanke and James(1991)]{BernankeJames1991}
Bernanke, B. and H.~James (1991).
\newblock The gold standard, deflation, and financial crisis in the great
  depression: An international comparison.
\newblock In R.~G. Hubbard (Ed.), \textit{Financial Markets and Financial
  Crises}, pp. 33--68. University of Chicago Press.

\bibitem[{Bullion Committee}(1810)]{BullionReport1810}
{Bullion Committee} (1810).
\newblock \textit{Report together with Minutes of Evidence, and Accounts, from
  the Select Committee on the High Price of Gold Bullion}.
\newblock House of Commons, London.

\bibitem[Bordo(1990)]{Bordo1990}
Bordo, M.~D. (1990).
\newblock The lender of last resort: Alternative views and historical
  experience.
\newblock \textit{Federal Reserve Bank of Richmond Economic Review} 76(1),
  18--29.

\bibitem[Eichengreen(1992)]{Eichengreen1992}
Eichengreen, B. (1992).
\newblock \textit{Golden Fetters: The Gold Standard and the Great Depression,
  1919--1939}.
\newblock Oxford University Press.

\bibitem[Eichengreen(1996)]{Eichengreen1996}
Eichengreen, B. (1996).
\newblock \textit{Globalizing Capital: A History of the International Monetary
  System}.
\newblock Princeton University Press.

\bibitem[Fetter(1965)]{Fetter1965}
Fetter, F.~W. (1965).
\newblock \textit{Development of British Monetary Orthodoxy, 1797--1875}.
\newblock Harvard University Press.

\bibitem[Fisher(1911)]{Fisher1911}
Fisher, I. (1911).
\newblock \textit{The Purchasing Power of Money: Its Determination and Relation
  to Credit, Interest and Crises}.
\newblock Macmillan.

\bibitem[Fisher(1920)]{Fisher1920}
Fisher, I. (1920).
\newblock \textit{Stabilizing the Dollar}.
\newblock Macmillan.

\bibitem[Fisher(1935)]{Fisher1935}
Fisher, I. (1935).
\newblock \textit{100\% Money}.
\newblock Adelphi.

\bibitem[Friedman(1960)]{Friedman1960}
Friedman, M. (1960).
\newblock \textit{A Program for Monetary Stability}.
\newblock Fordham University Press.

\bibitem[Friedman(1962)]{Friedman1962}
Friedman, M. (1962).
\newblock \textit{Capitalism and Freedom}.
\newblock University of Chicago Press.

\bibitem[Friedman(1969)]{Friedman1969}
Friedman, M. (1969).
\newblock The optimum quantity of money.
\newblock In \textit{The Optimum Quantity of Money and Other Essays}. Aldine.

\bibitem[Friedman(2003)]{Friedman2003}
Friedman, M. (2003).
\newblock The Fed's thermostat.
\newblock \textit{Wall Street Journal}, August~19.

\bibitem[Friedman and Schwartz(1963)]{FriedmanSchwartz1963}
Friedman, M. and A.~J. Schwartz (1963).
\newblock \textit{A Monetary History of the United States, 1867--1960}.
\newblock Princeton University Press.

\bibitem[Humphrey(2001)]{Humphrey2001}
Humphrey, T.~M. (2001).
\newblock The real bills doctrine.
\newblock In D.~Glasner (Ed.), \textit{Business Cycles and Depressions: An
  Encyclopedia}. Routledge.

\bibitem[Humphrey and Timberlake(2019)]{HumphreyTimberlake2019}
Humphrey, T.~M. and R.~H. Timberlake (2019).
\newblock \textit{Gold, the Real Bills Doctrine, and the Fed: Sources of
  Monetary Disorder, 1922--1938}.
\newblock Cato Institute Press.

\bibitem[Keynes(1913)]{Keynes1913}
Keynes, J.~M. (1913).
\newblock \textit{Indian Currency and Finance}.
\newblock Macmillan.

\bibitem[Keynes(1923)]{Keynes1923}
Keynes, J.~M. (1923).
\newblock \textit{A Tract on Monetary Reform}.
\newblock Macmillan.

\bibitem[Keynes(1925)]{Keynes1925}
Keynes, J.~M. (1925).
\newblock \textit{The Economic Consequences of Mr.~Churchill}.
\newblock Hogarth Press.

\bibitem[Keynes(1930)]{Keynes1930}
Keynes, J.~M. (1930).
\newblock \textit{A Treatise on Money}. 2 vols.
\newblock Macmillan.

\bibitem[Keynes(1936)]{Keynes1936}
Keynes, J.~M. (1936).
\newblock \textit{The General Theory of Employment, Interest and Money}.
\newblock Macmillan.

\bibitem[Keynes(1943)]{Keynes1943}
Keynes, J.~M. (1943).
\newblock Proposals for an international clearing union.
\newblock \textit{British Government White Paper}, Cmd.\ 6437.

\bibitem[Laidler(1984)]{Laidler1984}
Laidler, D. (1984).
\newblock Misconceptions about the real-bills doctrine: A comment on Sargent
  and Wallace.
\newblock \textit{Journal of Political Economy} 92(1), 149--155.

\bibitem[Lucas(1992)]{Lucas1992}
Lucas, R.~E. (1992).
\newblock On efficiency and distribution.
\newblock \textit{Economic Journal} 102(411), 233--247.

\bibitem[Mints(1945)]{Mints1945}
Mints, L.~W. (1945).
\newblock \textit{A History of Banking Theory in Great Britain and the United
  States}.
\newblock University of Chicago Press.

\bibitem[Mises(1949)]{Mises1949}
Mises, L.~von (1949).
\newblock \textit{Human Action: A Treatise on Economics}.
\newblock Yale University Press.

\bibitem[Nurkse(1944)]{Nurkse1944}
Nurkse, R. (1944).
\newblock \textit{International Currency Experience: Lessons of the Interwar
  Period}.
\newblock League of Nations.

\bibitem[Ricardo(1810)]{Ricardo1810}
Ricardo, D. (1810).
\newblock \textit{The High Price of Bullion, A Proof of the Depreciation of
  Bank Notes}.
\newblock John Murray.

\bibitem[Ricardo(1816)]{Ricardo1816}
Ricardo, D. (1816).
\newblock \textit{Proposals for an Economical and Secure Currency: With
  Observations on the Profits of the Bank of England}.
\newblock John Murray.

\bibitem[Salerno(1982)]{Salerno1982}
Salerno, J.~T. (1982).
\newblock The gold standard: An analysis of some recent proposals.
\newblock \textit{Cato Policy Analysis} 16.

\bibitem[Sargent(1982)]{Sargent1982unpleasant}
Sargent, T.~J. (1982).
\newblock The ends of four big inflations.
\newblock In R.~E. Hall (Ed.), \textit{Inflation: Causes and Effects}.
  University of Chicago Press.

\bibitem[Sargent(2011)]{Sargent2011}
Sargent, T.~J. (2011).
\newblock Evolution or revolution?
\newblock \textit{American Economic Review} 101(3), 1--27.

\bibitem[Sargent and Wallace(1982)]{SargentWallace1982}
Sargent, T.~J. and N.~Wallace (1982).
\newblock The real-bills doctrine versus the quantity theory: A
  reconsideration.
\newblock \textit{Journal of Political Economy} 90(6), 1212--1236.

\bibitem[Smith(1776)]{Smith1776}
Smith, A. (1776).
\newblock \textit{An Inquiry into the Nature and Causes of the Wealth of
  Nations}. 2 vols.
\newblock Strahan and Cadell.

\bibitem[Thornton(1802)]{Thornton1802}
Thornton, H. (1802).
\newblock \textit{An Enquiry into the Nature and Effects of the Paper Credit of
  Great Britain}.
\newblock Hatchard.

\bibitem[Triffin(1960)]{Triffin1960}
Triffin, R. (1960).
\newblock \textit{Gold and the Dollar Crisis: The Future of Convertibility}.
\newblock Yale University Press.

\bibitem[Volcker and Gyohten(1992)]{Volcker1992}
Volcker, P. and T.~Gyohten (1992).
\newblock \textit{Changing Fortunes: The World's Money and the Threat to
  American Leadership}.
\newblock Times Books.

\end{thebibliography}

\end{document}
